\section{The five-dimensional Artin gate}
\label{sec:ma-s5}

\paragraph{Status after the Fourier argument.}
MA-FQ-004 below proves $\beta=b+n$ under the exact hypothesis
$(\mathcal R_B)$ stated in this section. The bounds, positive-excess
signatures, and formal order witness developed here remain valid
intermediate deductions; their arithmetic positive-excess branch
is excluded by that later theorem. The hypothesis
$(\mathcal R_B)$ is not proved for every extension in the family.
The input $M$ supplies it and also supplies the additional Hecke
zero-location hypotheses needed for GRH of all irreducible factors.

\paragraph{Notation and the audited input.}
We retain precisely Calegari's family: the Galois group of $L/\mathbb Q$
is $S_5$, complex conjugation is of type $(12)(34)$, and $5$ is
unramified with Frobenius of the same type. Write
$E=L^{A_5}$, $F=L^{F_{20}}$, $H=EF$, and let $\eta$ be the character
of $E/\mathbb Q$. In the character-table notation of
GRH\_NIGHT MT-S506, put
\[
 V=\rho_4,\quad W=\rho_5,\quad X=\rho_6,\qquad
 B=\zeta_H/\zeta_F=L_F(s,\eta_F),\quad C=L(s,\eta),
 \quad A=L_f(s,\varpi).
\]
Here $\varpi$ is Calegari's specified cuspidal comparison representation,
not an arbitrary degree-five automorphic representation. The notation
$A$ in this section is an $L$-function; the Dirichlet GRH theorem is
denoted $\mathsf A$ when it is needed. All finite $L$-functions include
the ramified factors. Set
\[
 \begin{gathered}
 \beta=\operatorname{ord}_{1/2} B,\quad b=\operatorname{ord}_{1/2} C,
 \quad n=\operatorname{ord}_{1/2} A,\quad
 k=\beta-b-n,\\
 (\mathcal R_B):\ B(\sigma)\ne0\quad
 (0<\sigma<1,\ \sigma\ne1/2,\ \sigma\in\mathbb R).
 \end{gathered}
\]
We use the two all-prime identities
\[
 B=C L_W,\qquad \zeta_F=\zeta_{\mathbb Q}L_{W\eta}.
 \tag{MA-S5.1}
\]
They are GRH\_NIGHT MT-S506, equation (S506.6), and MT-S508,
equation (S508.2); their character calculations and induction conventions
were rechecked in this run.

\subsection*{MA-S5-001. What the earlier comparison actually proves}

\paragraph{Audited finding.}
The assertions MT-S504, MT-S505, MT-S507 and MT-S508 pass the checks
specified below, with their stated hypotheses. In particular, under
$(\mathcal R_B)$,
\[
 k\ge0,\qquad
 k=0\ \Longleftrightarrow\ \Omega\text{ is finite}
 \ \Longleftrightarrow\
 \varpi\text{ and }W\text{ have matching factors at every place}.
 \tag{MA-S5.2}
\]
The last condition also gives the all-place match for the quadratic
twist. It proves automorphy and holomorphy of the nontrivial irreducible
Artin factors in this family. GRH for all these factors requires the
additional zero-location input $M$, or separate corresponding Hecke
GRH hypotheses. The condition $(\mathcal R_B)$ alone is not such an input.

Calegari's initial comparison is equality on a set of primes of
Dirichlet density one. It is not equality at all but finitely many
places. Replacing the first assertion by the second would assume
the finiteness of $\Omega$ that is to be proved.

\paragraph{Complete check of the deductions used.}
The primary locators are Calegari (2013), Theorem 1.1 (author-hosted
publisher PDF p.~2), Theorem 4.7 (p.~12), Lemma 4.8 (pp.~13--14),
Landau Lemma 4.9 (p.~14), Lemma 4.10 (pp.~14--15), and the final
proof of Theorem 1.2 (p.~15). These are publisher-PDF page numbers,
not the arXiv-v1 page numbers used in MT-S504. The complete proofs
of the comparison lemmas, their intervening Landau argument, and
the final deduction were read. The construction of the cuspidal
representations remains Calegari's cited theorem.

There are also harmless sign/asymptotic slips in the printed proofs
of Lemma 4.6 and Theorem 4.7 (publisher-PDF p.~12). In the first,
the density-weighted coefficient is $-1/3$, so the logarithm of
the comparison quotient is
$-\frac13\log(1/(s-1))+O(1)$ as $s\downarrow1$.
The quotient is not just holomorphic but nonzero at $1$, since
each of its four factors has an exact simple pole there. Its
bounded logarithm supplies the contradiction with the corrected
asymptotic. In the second, the corrected identity is
\[
 \log\left|\frac{L(s,\Pi\times\Pi^\vee)}
                       {L(s,\operatorname{End}\xi)}\right|
       =8\sum_{p\in\Omega}p^{-s}+O(1).
\]
The left side is bounded for the same simple-pole reason.
Monotone convergence therefore gives
$\sum_{p\in\Omega}p^{-1}<\infty$, which is stronger than density
zero. The printed $o(|\log(s-1)|^{-1})$ does not follow from
boundedness; the valid density conclusion uses
$O(1)=o(\log(1/(s-1)))$. These corrections do not change any
comparison theorem used here.

For clarity, with $G$ denoting the archimedean factors, the quotient
that is actually used is
\[
 R(s)=\frac{G_{\varpi}(s)A(s)}{G_W(s)L_W(s)}
     =R_{\mathcal B}(s)D_\Omega(s),\qquad
 D_\Omega(s)=\prod_{p\in\Omega}
       \left(\frac{1+p^{-s}}{1-p^{-s}}\right)^2.
 \tag{MA-S5.3}
\]
For the convention $L_p=P_p(p^{-s})^{-1}$, the finite factor in
$R_{\mathcal B}$ is $P_{W,p}/P_{\varpi,p}$. The reversed displayed
finite-polynomial ratio in Calegari's Lemma 4.8 is the normalization
slip already corrected in MT-S508 (S508.9); it does not change the
real-interval nonvanishing of the finite factor. At $p\in\Omega$,
\[
 P_{\varpi,p}(T)=(1+T^2)(1-T)^3,\qquad
 P_{W,p}(T)=(1+T^2)(1-T)(1+T)^2.
\]
The finite and archimedean factor $R_{\mathcal B}$ is holomorphic
and nonzero at real $0<s<1$. If $\Omega$ is infinite, its positive
Euler series has an abscissa $\gamma\in(0,1)$ and a pole at
$\gamma$. The strict upper bound uses convergence at $1$, which
follows from the proof of Theorem 4.7, rather than merely from
density zero. The exclusion of $\gamma=0$ is the arbitrary-finite-subset
argument on publisher p.~14. Thus Lemmas 4.8--4.9 supply an actual
real pole, and Lemma 4.10 supplies exact comparison if $\Omega$
is finite.

Substitution of (MA-S5.1) gives
\[
 R(s)=\frac{G_{\varpi}(s)}{G_W(s)}\frac{A(s)C(s)}{B(s)},
 \qquad \operatorname{ord}_{1/2}R=-k.
\]
Under $(\mathcal R_B)$ every possible real pole is at $1/2$.
Infinite $\Omega$ therefore implies $k>0$; finite $\Omega$ gives
$R=1$ and $k=0$. These two exhaustive cases prove the central
equivalence and its inequality. This is a check of the proof of
MT-S508, not a new proof of its automorphic construction.

MT-S504 assumes the stronger input $M$ and $\operatorname{ord}_{1/2}L_W\le0$.
Its signed Brauer argument makes $L_W$ nonzero and holomorphic on
the other real points; a real pole of $R$ would then require a
positive central order of $L_W$. Consequently its implication is
valid. MT-S505 correctly uses that implication to exclude a negative
central order, then uses $\zeta_{\mathbb Q}(\sigma)<0$ on $(0,1)$
in the second identity (MA-S5.1). It does not infer nonvanishing
of $L_W$ from its holomorphy. The parity part concerns meromorphic
orders, for which the Laurent-series argument is valid.

For MT-S507 the whole displayed cone proof and its sharpness proof
were checked. Its premise $c,d\ge0$ is an extra holomorphy input.
The exhaustive finite calculation was re-executed after removing
the exploratory optimizer: $156$ subgroups, $19$ conjugacy classes,
and $43$ induced-linear multiplicity vectors were recovered. Every
vector was reconstructed as a nonnegative integer combination of
the twelve displayed generators. The four sharpness witnesses
were checked in integer arithmetic. The files in
\texttt{checks/s5/} contain that computation. These witnesses are
cone points, not number fields.

Finally, the half-plane and one-$B$ consequences in MT-S508 were
checked with their precise scope. Full GRH for $B$ gives GRH for
$A$ and $C$ and critical-line divisor support for $L_W$. It does
not by itself give GRH for $L_V,L_{V\eta},L_X,L_{W\eta}$.
This last distinction is needed when interpreting task 2's
phrase ``holomorphy and GRH for every irreducible factor.''

\subsection*{MA-S5-002. The residual excess is even}

\paragraph{Statement.}
In this family the three orders $\beta,b,n$ are even, independently
of $(\mathcal R_B)$. Consequently, under $(\mathcal R_B)$,
\[
        k\in 2\mathbb Z_{\ge0}.
\]
At this stage the hypothetical positive branch of MT-S508 is therefore
\[
             k\ge2\quad\text{and }k\text{ even},
 \tag{MA-S5.4}
\]
not the potentially odd excess allowed by its stated bound $k\ge1$.
The orders $n_\eta=\operatorname{ord}_{1/2}L_f(s,\varpi\otimes\eta)$
and $f=\operatorname{ord}_{1/2}\zeta_F$ are even as well.

\begin{proof}
The completed Dedekind zeta functions have functional equation with
sign $+1$. The completed quotients for $H/F$ and $E/\mathbb Q$
therefore give sign $+1$ for the quadratic Hecke functions $B$ and
$C$. Their gamma and conductor factors are holomorphic and nonzero
at $1/2$. If a nonzero meromorphic function satisfies
$\Lambda(1/2+t)=\Lambda(1/2-t)$, its leading Laurent term
$a t^r$ satisfies $a=(-1)^r a$, so $r$ is even. Applying this to
$B,C$, and $\zeta_F$ gives the first two parities and the assertion
about $f$.

At every unramified prime, the Satake multiset of $\varpi$ is
invariant under inversion. Away from $\Omega$ it is the multiset
of the real orthogonal representation $W$; on $\Omega$ it is
$\{1,1,1,i,-i\}$. Thus $\varpi$ and $\varpi^\vee$ agree at all
but finitely many places. Ordinary strong multiplicity one applies
to these two already automorphic representations, and gives
$\varpi\simeq\varpi^\vee$. Notice that this use of multiplicity
one is different from the unavailable all-but-finite comparison
of $\varpi$ with $W$.

A self-dual cuspidal representation of odd degree is of orthogonal
type. Arthur (2013), Theorems 1.4.1 and 1.5.3(a), gives this
orthogonal/symplectic classification, and Theorem 1.5.3(b), applied
also to the trivial representation of $\mathrm{GL}_1$, gives
$\epsilon(1/2,\varpi)=+1$. The exact epsilon-factor statement is
also stated by Arthur as Theorem 3(b) in his author-hosted
\emph{The Endoscopic Classification of Representations}, p.~8;
the specialization is explicitly recorded as Theorem 3.3.6 in
Dalal--Gerbelli-Gauthier, \texttt{arXiv:2410.01976v2}. The input
here is Arthur's classical theorem, not a new theorem from that
preprint. Since the completed standard $L$-function is holomorphic
at the centre and its gamma factor is nonzero there, the same
Laurent argument gives $n\in2\mathbb Z_{\ge0}$.

The quadratic twist is again self-dual cuspidal of odd degree, so
the same argument gives even $n_\eta$. Finally $k=\beta-b-n$ is
even, and MT-S508 supplies $k\ge0$ under $(\mathcal R_B)$.
No zero-location statement was used in any parity argument.
\end{proof}

\paragraph{Why MT-ART10--MT-ART12 do not close this gate.}
Those statements were read with their hypotheses: the integral
class and hermitian invariants in MT-ART11--MT-ART12 give
functional-equation parities, particularly on the symplectic
character group, and MT-ART10's forced zero is a quaternionic
example. The present $S_5$ irreducibles are real orthogonal and
have root number $+1$ (Calegari, p.~3). No negative sign is available
to force or rule out a positive even excess. The additional
automorphic root-number input above does eliminate the odd values
of $k$, but parity alone does not imply $k=0$.

\subsection*{MA-S5-003. A Rankin--Selberg bound for the excess}

\paragraph{Statement.}
Let $\Pi$ and $\xi$ be Calegari's four-dimensional automorphic
representation and finite-image representation of the central
extension of $S_5$ used in his section 4. Put
\[
 c=\operatorname{ord}_{1/2}L_V,\quad
 d=\operatorname{ord}_{1/2}L_{V\eta},\quad
 e=\operatorname{ord}_{1/2}L_X,
\]
and let $r_4,r_5$ be the central orders of the completed
Rankin--Selberg functions of $\Pi\times\Pi^\vee$ and
$\varpi\times\varpi^\vee$, respectively. Then
\[
 r_4=n+d+e-k,\qquad
 r_5=r_4+c+f.
 \tag{MA-S5.5}
\]
In particular, under $(\mathcal R_B)$,
\[
     0\le k\le n+d+e.
 \tag{MA-S5.6}
\]
If $A(1/2)L_{V\eta}(1/2)L_X(1/2)\ne0$, the comparison closes.
Combining with MA-S5-002, failure to close requires
\[
 2\le k\le n+d+e,\qquad k\in2\mathbb Z,
 \qquad n,d,e\in2\mathbb Z_{\ge0}.
 \tag{MA-S5.7}
\]
Thus it requires a central zero in at least one of these three
known entire automorphic functions. This is a necessary condition;
it is not an asserted nonvanishing theorem for them.

\begin{proof}
We first make the finite representation calculation explicit.
Scalar central elements act trivially on $\operatorname{End}\xi$,
so it is a representation of $S_5$. In the order of classes
$1,2A,2B,3A,4A,5A,6A$, its character is
\[
                    (16,0,0,1,0,1,3).
\]
For the odd classes, this is the squared absolute trace of the
projective matrices in Calegari's section 2, publisher-PDF p.~6:
the traces are $0,0,\sqrt3$ up to scalar. On $A_5$, it follows
from $\xi|_{N.A_5}=\operatorname{Sym}^3\varrho$. Indeed the
symmetric-cube traces at projective orders $2,3,5$ have absolute
squares $0,1,1$, respectively. For example, if the two eigenvalues
of a determinant-one lift are $e^{i\theta},e^{-i\theta}$, the
symmetric-cube trace is
$2\cos3\theta+2\cos\theta$; taking $\theta=\pi/2,\pi/3,\pi/5$
or $2\pi/5$ gives these squares. Comparing every entry with
the $S_5$ table proves
\[
          \operatorname{End}\xi=\mathbf1\oplus V\eta\oplus W\oplus X.
 \tag{MA-S5.8}
\]
An entrywise product and inner product in the same character table
gives
\[
 W\otimes W=\mathbf1\oplus V\oplus V\eta\oplus W\oplus W\eta\oplus X.
 \tag{MA-S5.9}
\]
Both identities, and the Euler-polynomial calculation below, have
an exact integer/cyclotomic certificate in \texttt{checks/s5/}.

Enlarge the finite bad-prime set $\mathcal B$ if needed. At a good
prime outside $\Omega$, the Satake multiset of $\Pi$ is the
multiset of $\xi$ up to a common sign; their endomorphism
multisets therefore agree. At $p\in\Omega$, up to a scalar root
of unity and replacing $i$ by $-i$, the two multisets are
\[
          S_\Pi=\{1,i,1,i\},\qquad S_\xi=\{1,i,-1,-i\}.
\]
These are the second and fourth exceptional rows of Calegari's
$4A$ table, publisher p.~11. For every positive integer $j$,
\[
 \left|\sum_{\alpha\in S_\Pi}\alpha^j\right|^2
 -\left|\sum_{\alpha\in S_\xi}\alpha^j\right|^2
 =\begin{cases}8,&j\text{ odd},\\0,&j\text{ even}.\end{cases}
 \tag{MA-S5.10}
\]
For odd $j$ the first squared trace is $8$ and the second is
$0$; for $j\equiv2\pmod4$ both are $0$; for $4\mid j$ both
are $16$. The scalar root of unity cancels in each endomorphism
trace. Expanding Euler logarithms in $\Re s>1$ gives
\[
 \log\frac{L_p(s,\Pi\times\Pi^\vee)}{L_p(s,\operatorname{End}\xi)}
 =8\sum_{j\text{ odd}}\frac{p^{-js}}j
 =2\log\left(\frac{1+p^{-s}}{1-p^{-s}}\right)^2.
\]
Multiplying at good primes, then completing, yields a meromorphic
identity
\[
 \frac{\Lambda(s,\Pi\times\Pi^\vee)}
      {\Lambda(s,\operatorname{End}\xi)}
          =U_4(s)D_\Omega(s)^2,
 \tag{MA-S5.11}
\]
where $U_4$ consists of finitely many local factors, gamma factors,
and conductor powers and is holomorphic and nonzero at $1/2$.

Here is the check of that last, potentially important assertion.
The base change of $\Pi$ to $E$ has the finite-image local
parameters of $\operatorname{Sym}^3\varrho$. Restriction of a
Weil--Deligne parameter to the finite-index local Weil group
over $E$ has zero monodromy and finite image; consequently the
original parameter has zero monodromy and finite image as well.
Its image is a union of finitely many translates of the finite
restricted image. Thus its endomorphism parameter has eigenvalues
that are roots of unity, including at bad primes. No polynomial
$\det(1-p^{-s}\operatorname{Frob})$ can vanish at $s=1/2$.
The Artin parameter has the same property. Since $E$ is real,
base change is split at infinity, where the parameters are also
finite-image parameters; the gamma factors are products of
$\Gamma_{\mathbb R}(s)$ and $\Gamma_{\mathbb R}(s+1)$ and are
nonzero and finite at the centre. Conductor powers never vanish.
This proves the assertion about $U_4$ without an all-place
matching assumption.

By (MA-S5.3), $\operatorname{ord}_{1/2}D_\Omega=-k$.
Artin multiplicativity in (MA-S5.8), together with
$\zeta_{\mathbb Q}(1/2)\ne0$, gives
\[
 \operatorname{ord}_{1/2}\Lambda(s,\operatorname{End}\xi)=d+x+e,
 \qquad x=\operatorname{ord}_{1/2}L_W=\beta-b=n+k.
\]
Taking orders in (MA-S5.11) therefore proves
$r_4-(d+x+e)=-2k$, or the first identity in (MA-S5.5).

For the degree-five Rankin--Selberg comparison, the odd-power
traces at an exceptional prime are $3$ and $-1$, and the even-power
traces agree. Their squared-trace difference is again $8$ on odd
powers and $0$ on even powers. The same calculation, now with
(MA-S5.9), gives
\[
 r_5=c+d+x+y+e-2k,
 \qquad y=\operatorname{ord}_{1/2}L_{W\eta}=f.
\]
Subtracting the formula for $r_4$ proves $r_5-r_4=c+f$.
The bad-factor argument is identical, using the finite-image
base change of $\varpi$.

The completed Rankin--Selberg function of two unitary cuspidal
representations is holomorphic at $1/2$: its possible poles lie
on $\Re s=0$ or $1$. An exact usable locator is Cogdell (2004),
\emph{Lectures on $L$-functions, Converse Theorems, and Functoriality
for $\mathrm{GL}_n$}, Theorem 9.2(i), printed p.~74 (PDF p.~78),
which states the Jacquet--Shalika analytic result. Thus $r_4\ge0$.
Calegari's Theorem 1.1 gives $d,e\ge0$, and $n\ge0$ because
$A$ is entire. Under $(\mathcal R_B)$, MT-S508 gives $k\ge0$.
These prove (MA-S5.6). If the stated three central values are
nonzero, the upper bound is zero. Finally $d,e$ are even by the
root-number-$+1$ functional equations of the real $S_5$
representations, and MA-S5-002 proves the remaining parities.
\end{proof}

\subsection*{MA-S5-006. Symmetric-square holomorphy removes the comparison order}

\paragraph{Statement.}
With the same notation, the sharper bound is
\[
                 0\le k\le d+e
       \qquad\text{under }(\mathcal R_B).
 \tag{MA-S5.16}
\]
Consequently $L_{V\eta}(1/2)L_X(1/2)\ne0$ suffices to close the
comparison, regardless of the central order of $A$. A hypothetical
positive branch must satisfy
\[
       2\le k\le d+e,\qquad k,d,e\in2\mathbb Z_{\ge0}.
 \tag{MA-S5.17}
\]

\begin{proof}
Let $\nu=\psi^{-1}$ be Calegari's finite-order similitude character.
The exact inputs are his Lemma 4.2 and the following displayed
decomposition (publisher-PDF pp.~9--10), and Lemma 4.6
(publisher-PDF p.~12):
\[
 \Pi^\vee=\Pi\otimes\nu^{-1},\qquad
       (\wedge^2\Pi)\otimes\nu^{-1}=\mathbf1\boxplus\varpi.
\]
Choose a finite set $S$ containing the archimedean places, the
places above $2$, and all ramified places in these identities.
At every unramified place, the decomposition of the tensor square
into symmetric and exterior squares gives
\[
 L^S(s,\Pi\times\Pi^\vee)
  =L^S(s,\Pi,\operatorname{Sym}^2\otimes\nu^{-1})
       \zeta_{\mathbb Q}^S(s)A^S(s).
\]
This identity first holds in an absolute-convergence half-plane
and then meromorphically. The omitted factors are holomorphic
and nonzero at $1/2$, by the finite-image parameter check in
MA-S5-003. The character $\nu$ has finite order, and the central
character of $\Pi$ is unitary: its base change has finite-order
central character, so a power of it has finite order already
over $\mathbb Q$. Thus Takeda (2015), Theorem 7.1, applies to
the twisted symmetric-square factor. It is holomorphic at $1/2$.

Write $q$ for its central order. Since $\zeta_{\mathbb Q}(1/2)\ne0$,
the displayed identity gives $r_4=q+n$. MA-S5-003 gives
$r_4=n+d+e-k$. Hence
\[
                         q=d+e-k\ge0.
\]
Combining this with MT-S508 proves (MA-S5.16); MA-S5-002 and the
Artin parities already checked give (MA-S5.17). No nonvanishing
of the symmetric-square factor has been asserted or used.
\end{proof}

\paragraph{Source-check boundary.}
Takeda's result is the published theorem in
\emph{Mathematische Zeitschrift} 281 (2015), 103--157, not an
additional conjectural preprint assumption. The exact primary
locator read here is Theorem 7.1 and its complete terminal proof,
\texttt{arXiv:1506.04791v1}, pp.~49--50; that proof reduces the
claim to its Theorem 6.2 and the earlier integral representation.
The latter classical upstream results were not independently
re-proved here. The new deduction above and its local
normalizations were checked in full. The representation-ring
form of the same deduction is
$(\operatorname{Sym}^2\xi)\otimes\nu^{-1}=V\eta\oplus X$,
obtained by subtracting
$(\wedge^2\xi)\otimes\nu^{-1}=\mathbf1\oplus W$
from (MA-S5.8). Thus this is a symmetric-power statement in the
same Artin-motive $\lambda$-ring as MA-S5-003.

\subsection*{MA-S5-004. A stronger exceptional-prime signature}

\paragraph{Statement.}
Assume $(\mathcal R_B)$ and $k>0$. Besides MT-S508's Abelian law,
one has the ordinary weighted asymptotic
\[
  \sum_{\substack{p\le X\\p\in\Omega}}
           \frac{\log p}{\sqrt p}
            \sim \frac{k}{4}\log X,
 \tag{MA-S5.12}
\]
and, writing $\pi_\Omega(X)=\#\{p\le X:p\in\Omega\}$,
\[
                 \lim_{X\to\infty}
                  \frac{\log \pi_\Omega(X)}{\log X}=\frac12.
 \tag{MA-S5.13}
\]
This is an exponent statement for the unweighted counting
function, not an asymptotic $\pi_\Omega(X)\sim c\sqrt X/\log X$.
In particular any upper bound $\pi_\Omega(X)=O(X^{1/2-\delta})$
with fixed $\delta>0$ would close the comparison.

\begin{proof}
The proof of MT-S508 gives absolute convergence of the Euler
logarithm in $\Re s>1/2$ and
\[
 \log D_\Omega(s)=4\sum_{p\in\Omega}p^{-s}+E(s),
\]
where $E$ is holomorphic in $\Re s>1/3$; its proof is the
locally uniform comparison with $\sum_p p^{-3\Re s}$. Since
$D_\Omega$ has a pole of order $k$ at $1/2$, logarithmic
differentiation along $s=1/2+\epsilon$ gives
\[
 \sum_{p\in\Omega}\frac{\log p}{p^{1/2+\epsilon}}
                  =\frac{k}{4\epsilon}+O(1).
 \tag{MA-S5.14}
\]
We give the positive-measure Tauberian step explicitly, rather
than use an unstated strengthening of the Abelian law.
Let $\nu$ be the positive measure on $[0,\infty)$ with mass
$(\log p)/\sqrt p$ at $\log p$ for $p\in\Omega$, and put
$a=k/4$. Its Laplace transform $F(t)$ satisfies
$F(t)\sim a/t$ as $t\downarrow0$. Define positive measures
$\nu_T$ by $\nu_T(J)=T^{-1}\nu(TJ)$. Then for every $v>0$,
\[
       \int_0^\infty e^{-v x}\,d\nu_T(x)
            =T^{-1}F(v/T)\longrightarrow a/v.
\]
The finite measures $d\lambda_T(x)=e^{-x}d\nu_T(x)$ have
total mass tending to $a$. Their tails satisfy, for $R>0$,
\[
 \lambda_T([R,\infty))
   \le e^{-R/2}\int e^{-x/2}\,d\nu_T(x)
   =e^{-R/2}T^{-1}F(1/(2T)),
\]
which is uniformly small for large $T$ as $R\to\infty$.
Thus they are tight. Every weakly convergent subsequence has
Laplace transforms $a/(v+1)$, $v\ge0$, and hence equals
$a e^{-x}dx$. For completeness, uniqueness follows by pushing
the measures forward under $x\mapsto e^{-x}$ to $[0,1]$:
equality of the Laplace transforms at nonnegative integers is
equality of all polynomial moments, and polynomials are dense
in the continuous functions on $[0,1]$. Therefore all of
$\lambda_T$ converge weakly to $a e^{-x}dx$.
Multiplying by $e^x$ on the compact interval $[0,1]$ and
approximating its indicator from above and below by continuous
functions gives
\[
           T^{-1}\nu([0,T])=\nu_T([0,1])\longrightarrow a.
\]
There are no atoms at the endpoints in the limiting measure,
so this last approximation introduces no boundary term. With
$T=\log X$, it proves (MA-S5.12).

For each fixed $0<u<1$, subtraction of (MA-S5.12) at $X$ and
$X^u$ gives
\[
 \sum_{\substack{X^u<p\le X\\p\in\Omega}}
       \frac{\log p}{\sqrt p}
              \sim a(1-u)\log X.
\]
For large $X$, the decreasing function $(\log t)/\sqrt t$
is at most $u(\log X)X^{-u/2}$ on this interval. Consequently
\[
 \pi_\Omega(X)\ge
     \left(\frac{a(1-u)}u+o(1)\right)X^{u/2}.
\]
This proves that the lower limit in (MA-S5.13) is at least
$u/2$, and then at least $1/2$ by letting $u\uparrow1$.
Conversely, absolute convergence of
$\sum_{p\in\Omega}p^{-\sigma}$ for every $\sigma>1/2$
gives $\pi_\Omega(X)\le X^\sigma\sum_{p\in\Omega}p^{-\sigma}$.
The upper limit is at most $\sigma$, hence at most $1/2$.
This proves the limit and the final criterion.
\end{proof}

\subsection*{MA-S5-005. The checked central constraints still permit positive excess}

\paragraph{Statement and its scope.}
The parity, factorization, complete monomial-cone inequalities,
and new Rankin--Selberg and symmetric-square central inequalities
do not, as a system
of constraints on integer orders, force $k=0$. This statement
does not construct a number field and does not establish that
$k>0$ occurs arithmetically.

\begin{proof}
For the coordinates of MT-S506 take
\[
 (a,b,c,d,x,y,e)=(0,0,0,2,2,0,0),
 \qquad n=0,\quad n_\eta=2,\quad\beta=2,\quad f=0.
 \tag{MA-S5.15}
\]
Then $k=\beta-b-n=2=x-n=n_\eta-f$. Every coordinate that is
required to be nonnegative is nonnegative and even. The seven
nontrivial generators in MT-S507 are
\[
 (x+y,b+x,a+y,c+x+e,d+y+e,a+c+x,b+d+y)
                       =(2,2,0,2,2,2,2).
\]
Hence all $43$ monomial tests are nonnegative. The new formulas
give $r_4=n+d+e-k=0$ and $r_5=r_4+c+f=0$. All relevant
factorization orders agree, and the symmetric-square order is
$q=d+e-k=0$. This verifies the claimed
compatibility in integer arithmetic.

It can also be realized by even meromorphic germs, solely to
check the local functional-equation and multiplicative relations.
With $t=s-1/2$, set $L_W=t^2$, $L_{V\eta}=t^2$,
$A_\eta=t^2$, $B=t^2$, and set $A,C,L_V,L_X,L_{W\eta}$ and
the zeta-base germ equal to $1$. Put $D=t^{-2}$ and both
Rankin--Selberg germs equal to $1$. The two tensor-product
identities then have the required orders and every germ is even.
These germs have no prescribed arithmetic Euler factors. They
are a witness to the insufficiency of the listed central
constraints, not a counterexample to the Artin conjecture or to
the requested arithmetic equality.
\end{proof}

\subsection*{MA-S5-007. The additional input needed by the multiplicity-one route}

\paragraph{A sufficient extra hypothesis, not an input used above.}
Suppose that there exists an isobaric sum $\Theta$ of unitary cuspidal
automorphic representations of $\mathrm{GL}_{10}/\mathbb Q$ whose
unramified local factors agree, outside a finite set, with
$L(s,\Pi,\operatorname{Sym}^2\otimes\nu^{-1})$.
Then $\Omega$ is finite and $k=0$. This is a conditional reduction;
Takeda's holomorphy theorem does not assert this automorphic
realization, and no such realization is supplied in this run.

\begin{proof}
Let $\Sigma=\pi_{V\eta}\boxplus\pi_X$, using the two cuspidal
transfers already supplied by Calegari. The identity at the end of
MA-S5-006 gives matching local factors for $\Theta$ and $\Sigma$
outside $\Omega$ and a finite set $S$. Their Satake parameters at
these good primes have modulus one. Write their cuspidal
decompositions with respect to their combined set of distinct
unitary cuspidal constituents $\tau_i$ as multiplicities
$m_{1,i},m_{2,i}\ge0$. Consider the quotient of partial functions
\[
 Q(s)=\frac{L^S(s,\Theta\times\Theta^\vee)
                 L^S(s,\Sigma\times\Sigma^\vee)}
                {L^S(s,\Theta\times\Sigma^\vee)
                 L^S(s,\Sigma\times\Theta^\vee)}.
\]
At a good prime its Euler logarithm is
$\sum_{j\ge1}|\operatorname{tr}T_{\Theta,p}^j-
\operatorname{tr}T_{\Sigma,p}^j|^2p^{-js}/j$.
It is zero outside $\Omega$. On $\Omega$ the traces are bounded
in absolute value by $10$, so $\sum_{p\in\Omega}p^{-1}<\infty$
from MA-S5-001 proves that $\log Q(\sigma)$ has a finite limit
as $\sigma\downarrow1$. Thus $Q$ has neither a zero nor a pole
at $1$.

For each pair of cuspidal constituents the partial
Rankin--Selberg function is nonzero at $1$ and has a simple pole
there precisely when the two constituents agree. Here the
simple-pole assertion is Cogdell (2004), Theorems 9.1--9.2 and the
partial-factor check on printed pp.~74--75; the nonvanishing is
Shahidi (1980), the theorem on p.~462. Multiplicativity therefore
gives
\[
 -\operatorname{ord}_{s=1}Q(s)
              =\sum_i(m_{1,i}-m_{2,i})^2.
\]
The left side is zero, so all multiplicities agree and
$\Theta=\Sigma$ as isobaric representations. At an exceptional
good prime their proposed local factors would have quotient
$((1+p^{-s})/(1-p^{-s}))^2$, which is not identically one.
Consequently all exceptional primes lie in the finite omitted
set. Calegari's Lemma 4.10 gives $A=L_W$, and (MA-S5.1)
then gives $\beta=b+n$, so $k=0$.
\end{proof}

\paragraph{Why a recent analytic multiplicity-one theorem does not provide it.}
Let $u_p,v_p$ be the prime coefficients of the twisted
symmetric-square function and $L_{V\eta}L_X$, respectively.
The exact local comparison gives $u_p-v_p=4$ on $\Omega$ and
$0$ outside it, apart from finitely many primes. Farmer, Pitale,
Ryan, and Schmidt (2026), Theorem 2.3, requires the distance
condition (2.20) in their revised primary preprint, pp.~11--12:
in particular, the series
\[
       \sum_p\frac{p|u_p-v_p|^2\log p}{p^\sigma}
 \tag{MA-S5.18}
\]
must converge for every $\sigma>1$. If $(\mathcal R_B)$ and
$k>0$ hold, MA-S5-004 gives
$\sum_{p\in\Omega,p\le X}(\log p)/\sqrt p
\sim(k/4)\log X$. Hence (MA-S5.18) diverges already at
$\sigma=3/2$, and by positivity also at every
$1<\sigma\le3/2$. Thus the theorem's hypothesis actually fails
under that hypothetical branch. Its full proof is not needed for this
hypothesis check and is not claimed independently audited.
This excludes that particular proposed application; it proves
neither the impossibility of a different analytic proof nor the
arithmetic existence of positive excess.

\paragraph{The successive reductions and the completed central argument.}
The Fourier argument in MA-FQ-001--004 proves $k=0$ under
$(\mathcal R_B)$. Before that argument, the checked necessary
conditions for a positive excess were (MA-S5.17), together with
(MA-S5.12)--(MA-S5.13).
The root-number route eliminates odd excess, the factorization
route isolates $k$ exactly, and the Rankin--Selberg route forces
additional central zeros; the symmetric-square route improves
the gate further to (MA-S5.17). Ordinary strong multiplicity one
cannot start from density-one comparison. The finite monomial
cone cannot supply the missing noncoincidence or nonvanishing:
(MA-S5.15) satisfies it even after the new inequalities.
Treating the meromorphic degree-zero quotient as a member of
the ordinary Selberg class would also assume a missing
holomorphy/finitely-many-poles condition. Finally, producing a
global finite-image Galois representation from the bounded local
parameters of the nonregular automorphic $\varpi$ is a separate
unproved realization step in this argument; the local parameters
alone do not furnish it. Nonvanishing of $L_{V\eta}L_X$, an
exceptional-prime estimate stronger than exponent $1/2$, and
the extra functoriality in MA-S5-007 were possible sufficient
inputs at that stage; none is needed in the completed proof.
That proof uses the exact equality of the gamma kernels,
finite-total-variation removal of the bad Euler terms, atomic
positive definiteness of the integer divisor, and the coset-ring
theorem. It excludes an infinite exceptional-prime set without
proving any of those additional central-value or functoriality
hypotheses. The ordinary real-axis hypothesis $(\mathcal R_B)$
remains explicit.

\paragraph{Grothendieck placement.}
Let $S$ contain the ramified primes and work over
$U=\operatorname{Spec}\mathbb Z[1/S]$. The $S_5$-cover gives
finite-monodromy lisse sheaves and their rational Artin motives.
SGA~1, Expos\'e~V, \S7 identifies finite \'etale covers with finite
continuous fundamental-group sets. Deligne (1979),
\emph{Valeurs de fonctions $L$ et p\'eriodes d'int\'egrales},
Definition 6.1 and \S\S6.2--6.3, identifies their Karoubian
rational linearization with Artin motives, whose realizations
have Hodge type $(0,0)$ and weight zero; equation (0.5.3) identifies
restriction of scalars with induction. Thus (MA-S5.1) is the
finite-pushforward decomposition $\operatorname{Ind}_{F_{20}}^{S_5}
\eta_F=\eta\oplus W$, while (MA-S5.8)--(MA-S5.9) are tensor
identities in the representation $K_0$ ring. Their exterior and
symmetric powers belong to the representation $\lambda$-ring
formalism, as placed in GRH\_NIGHT MT-ART11 and AG-LAMBDA
(SGA~6, Expos\'e~V, \S\S2,7, especially (7.1.1)).

More explicitly, MA-S5-001 concerns the induced Artin-motive
decomposition and the effectivity of its analytic divisor;
MA-S5-002 concerns the parity of that weight-zero motive's
functional equation together with the automorphic comparison.
MA-S5-003 concerns tensor products and the descended
endomorphism motive, and MA-S5-006 concerns its symmetric-power
summand. MA-S5-004 concerns the Frobenius places at which the
automorphic Satake data fail to equal the motive's Euler data;
its weighted counting measure is analytic and is not asserted
to be the dimension function of a cohomology group.
MA-S5-005 is only a compatible assignment of analytic orders
to the representation-ring identities, with no realization by
an Artin motive or finite \'etale cover. MA-S5-007 separates the
known symmetric-power Euler data of this Artin motive from the
extra automorphic realization needed by its conditional route.
These distinctions use
the same exact SGA~1 and Deligne locators just given.

The parity assertion is about the functional equation of a
weight-zero Artin motive and its known automorphic comparison;
the Rankin--Selberg bound is an inequality between analytic
central multiplicities. None of these constructions identifies
those multiplicities with dimensions of a positive faithful
arithmetic cohomology. They refine the Artin-divisor effectivity
gate of GRH\_NIGHT Q-ARTIN-POS and MT-S508. In the inherited
cote~1 trace/determinant programme (ALPOGE\_BRIDGES BR-TR-01,
BR-TR-02; sheets 79--80 and 111--112), they constrain a candidate
spectral divisor but do not supply the absent positivity theorem.
The latter archive locators are programme cross-references from
the supplied runs, not a claim that those scans were independently
re-read in this S5 audit.

\paragraph{Bibliographic locators used in this section.}
Calegari, F. (2013). \emph{The Artin conjecture for some
$S_5$-extensions}. \emph{Mathematische Annalen, 356}, 191--207.
\url{https://doi.org/10.1007/s00208-012-0839-4}.
Author-hosted publisher PDF:
\url{https://www.math.uchicago.edu/~fcale/papers/S5.pdf}.
The character table is on PDF p.~4; the projective $\xi$ table
is on p.~6; the exceptional $4A$ Satake table is on p.~11.

Arthur, J. (2013). \emph{The endoscopic classification of
representations: Orthogonal and symplectic groups} (Colloquium
Publications, Vol.~61). American Mathematical Society.
Theorems 1.4.1 and 1.5.3. The author's equivalent root-number
statement is Theorem 3(b), p.~8, at
\url{https://www.claymath.org/library/cw/arthur/pdf/arthur-endoscopic-tifr.pdf}.
The older overview's provisional account of the then unfinished
twisted trace formula is not used to replace the published theorem.

Cogdell, J. W. (2004). \emph{Lectures on $L$-functions, converse
theorems, and functoriality for $\mathrm{GL}_n$}. In
\emph{Lectures on automorphic $L$-functions} (Fields Institute
Monographs, Vol.~20). American Mathematical Society.
Author-hosted notes, Theorems 9.2 and 9.3, printed pp.~74--75:
\url{https://people.math.osu.edu/cogdell.1/fields-www.pdf}.

Takeda, S. (2015). On a certain metaplectic Eisenstein series and
the twisted symmetric square $L$-function.
\emph{Mathematische Zeitschrift, 281}, 103--157.
\url{https://doi.org/10.1007/s00209-015-1476-x}.
Theorem 7.1, author preprint pp.~49--50:
\url{https://arxiv.org/abs/1506.04791}.

Shahidi, F. (1980). On nonvanishing of $L$-functions.
\emph{Bulletin of the American Mathematical Society, 2}(3),
462--464. \url{https://doi.org/10.1090/S0273-0979-1980-14769-2}.
The theorem on p.~462 and its full printed outline on
pp.~462--463 were read; the outline's upstream Eisenstein-series
and Whittaker results remain classical inputs.

Farmer, D. W., Pitale, A., Ryan, N. C., \& Schmidt, R. (2026).
Multiplicity one for analytic $L$-functions and applications.
\emph{Mathematical Proceedings of the Cambridge Philosophical
Society, 181}(1), 885--911.
\url{https://doi.org/10.1017/S0305004126102011}.
The exact hypothesis checked is (2.20) and Theorem 2.3,
\texttt{arXiv:1305.3972v2}, pp.~11--12:
\url{https://arxiv.org/abs/1305.3972}.

Deligne, P. (1979). Valeurs de fonctions $L$ et p\'eriodes
d'int\'egrales. In \emph{Automorphic forms, representations and
$L$-functions} (Proceedings of Symposia in Pure Mathematics,
Vol.~33, Part~2, pp.~313--346). American Mathematical Society.
\S0.5.3 and \S\S6.1--6.3; Ghate--Beineke translation, PDF
pp.~4 and 20:
\url{https://mathweb.tifr.res.in/~eghate/Deligne.pdf}.
